% !TeX encoding = UTF-8
% !TeX root = ../main.tex

\section{整体算法流程}

\begin{frame}{整体算法流程与信号处理关联}
  \centering
  \small
  \begin{tikzpicture}[
    node distance=0.55cm and 0.5cm,
    box/.style={draw, rounded corners=2pt, minimum width=0.85cm, minimum height=0.55cm,
                align=center, font=\footnotesize},
    pre/.style={box, fill=blue!12},
    feat/.style={box, fill=green!12},
    cls/.style={box, fill=orange!15},
    output/.style={box, fill=red!12},
    arr/.style={->, >=stealth, thick},
  ]

  % --- Row 1: Preprocessing ---
  \node[pre] (raw) {Raw\\MP3};
  \node[pre, right=of raw] (ed) {端点\\检测};
  \node[pre, right=of ed] (pe) {预加重\\$y[n]=x[n]-0.97x[n-1]$};
  \node[pre, right=of pe] (frame) {分帧加窗\\25ms Ham, 10ms hop};

  % --- Row 2: Feature Extraction ---
  \node[feat, right=of frame] (fft) {FFT\\$|\cdot|^2$};
  \node[feat, right=of fft] (mel) {Mel 滤波\\26 个三角};
  \node[feat, right=of mel] (logdct) {Log + DCT\\13 维};
  \node[feat, right=of logdct] (mfcc) {MFCC\\+ $\Delta$/$\Delta^2$\\39 维};

  % --- Row 3: Classification ---
  \node[cls, below=1.0cm of mfcc] (clf) {KNN / GMM / CNN};
  \node[output, right=of clf] (pred) {预测\\0--9};

  % --- Arrows ---
  \draw[arr] (raw) -- (ed);
  \draw[arr] (ed) -- (pe);
  \draw[arr] (pe) -- (frame);
  \draw[arr] (frame) -- (fft);
  \draw[arr] (fft) -- (mel);
  \draw[arr] (mel) -- (logdct);
  \draw[arr] (logdct) -- (mfcc);
  \draw[arr] (mfcc.south) -- ++(0,-0.35) -| (clf.north);
  \draw[arr] (clf) -- (pred);

  % --- Labels ---
  \node[font=\footnotesize\bfseries, above=0.15cm of raw.north, anchor=south] {预处理};
  \node[font=\footnotesize\bfseries, above=0.15cm of fft.north, anchor=south] {特征提取};
  \node[font=\footnotesize\bfseries, below=0.15cm of clf.south, anchor=north] {分类器};

  \end{tikzpicture}

  \medskip
  \begin{columns}[T]
    \column{0.48\textwidth}
    \begin{block}{两条技术路线}
      \footnotesize
      \begin{itemize}
        \item \textbf{传统方法}：手工特征 + KNN（统计）/ GMM（序列）
        \item \textbf{深度方法}：MFCC 特征图 + 卷积神经网络（TinyKeywordCNN）
      \end{itemize}
    \end{block}

    \column{0.48\textwidth}
    \begin{block}{信号与系统核心关联}
      \footnotesize
      \begin{itemize}
        \item FFT：时域到频域的桥梁
        \item STFT：非平稳信号的时频联合分析
        \item Mel 滤波器组：模拟人耳非线性频率感知
        \item DCT：频域去相关与能量压缩
        \item 预加重/分帧加窗：FIR 滤波与短时平稳假设
      \end{itemize}
    \end{block}
  \end{columns}
\end{frame}
